% Регулятор переменного напряжения на встречно-параллельных тиристорах.
% Фазовое управление: VS1 проводит положительную полуволну,
% VS2 — отрицательную; нагрузка R_н.
\documentclass[border=4pt]{standalone}
\usepackage{amsmath}
\usepackage{fontspec}
\setmainfont{Calibri}
\usepackage[european, RPvoltages]{circuitikz}
\definecolor{ink}{RGB}{31,42,51}
\definecolor{grp}{RGB}{46,110,140}
\ctikzset{bipoles/length=1.0cm}
\ctikzset{bipoles/thickness=0.9}
\ctikzset{diodes/fill=ink}
\ctikzset{inductor=american}

\begin{document}
\begin{tikzpicture}[line width=1.0pt, color=ink, font=\small]
  % --- источник сети (синус)
  \draw (0,3.4) to[sV, l=$u_1$] (0,0);
  % --- верхняя шина до пары тиристоров
  \draw (0,3.4) -- (1.8,3.4) node[circ]{};
  % --- встречно-параллельные тиристоры между узлами P1 и P2
  \draw (1.8,3.4) to[Ty, name=VS1] (4.2,3.4) node[circ]{};
  \draw (1.8,3.4) -- (1.8,2.0);
  \draw (4.2,3.4) -- (4.2,2.0);
  \draw (4.2,2.0) to[Ty, name=VS2] (1.8,2.0);
  % --- подписи тиристоров (затворы УЭ — частью символа)
  \node[anchor=south] at (3.0,3.85) {VS1};
  \node[anchor=north] at (3.1,1.55) {VS2};
  % --- верхняя шина от P2 к нагрузке
  \draw (4.2,3.4) -- (6.0,3.4);
  % --- нагрузка
  \draw (6.0,3.4) to[R, l_=$R_{\text{н}}$] (6.0,0);
  % --- нижняя шина
  \draw (6.0,0) -- (0,0);
\end{tikzpicture}
\end{document}
